\documentclass[12pt,oneside]{book}

% ============================================================
% METADATOS
% ============================================================
\newcommand{\universidad}{Universidad de Buenos Aires (UBA)}
\newcommand{\facultad}{Facultad de Derecho}
\newcommand{\carrera}{Carrera de Abogacía}
\newcommand{\materia}{Práctico Profesional}
\newcommand{\titulo}{Derecho procesal civil en la práctica: negligencia, caducidad, costas e incidencias}
\newcommand{\autor}{Isaac Edgar Camacho Ocampo}
\newcommand{\anio}{Octubre de 2026}

% ============================================================
% PAQUETES
% ============================================================
\usepackage[utf8]{inputenc}
\usepackage[T1]{fontenc}
\usepackage[spanish,es-nodecimaldot]{babel}

\usepackage[
    a4paper,
    top=2.5cm,
    bottom=2.5cm,
    left=3cm,
    right=2.5cm,
    headheight=15pt
]{geometry}

\usepackage{amsmath}
\usepackage{amssymb}
\usepackage{mathtools}

\usepackage{graphicx}
\usepackage{float}
\usepackage{caption}
\usepackage{subcaption}

\usepackage{booktabs}
\usepackage{array}
\usepackage{longtable}
\usepackage{tabularx}

\usepackage{enumitem}

\usepackage{xcolor}
\usepackage[most]{tcolorbox}

\usepackage{fancyhdr}
\usepackage{ragged2e}

\usepackage[
    colorlinks=true,
    linkcolor=blue!50!black,
    citecolor=green!40!black,
    urlcolor=blue!60!black,
    pdftitle={\titulo},
    pdfauthor={\autor},
    pdfsubject={Apuntes universitarios},
    bookmarks=true
]{hyperref}

% ============================================================
% CONFIGURACIÓN
% ============================================================
\graphicspath{
    {./img/}
    {./graficos/}
    {./figuras/}
}

\setcounter{tocdepth}{2}

\setlength{\parindent}{0pt}
\setlength{\parskip}{0.5em}

\setlist[itemize]{
    topsep=0.3em,
    itemsep=0.2em
}

\setlist[enumerate]{
    topsep=0.3em,
    itemsep=0.2em
}

\clubpenalty=10000
\widowpenalty=10000
\emergencystretch=2em

\renewcommand{\arraystretch}{1.2}
\newcolumntype{L}[1]{>{\raggedright\arraybackslash}p{#1}}

% ============================================================
% COMANDOS
% ============================================================
\newcommand{\abs}[1]{\left|#1\right|}
\newcommand{\norm}[1]{\left\lVert#1\right\rVert}
\newcommand{\vect}[1]{\mathbf{#1}}
\newcommand{\deriv}[2]{\frac{d#1}{d#2}}
\newcommand{\pderiv}[2]{\frac{\partial#1}{\partial#2}}

% ============================================================
% ENTORNOS / CAJAS
% ============================================================
\newtcolorbox{definition}[1]{
    enhanced,
    breakable,
    title={Definición: #1},
    fonttitle=\bfseries,
    colback=blue!3,
    colframe=blue!50!black,
    boxrule=0.8pt,
    arc=2mm
}

\newtcolorbox{important}{
    enhanced,
    breakable,
    title={Importante},
    fonttitle=\bfseries,
    colback=yellow!8,
    colframe=orange!70!black,
    boxrule=0.8pt,
    arc=2mm
}

\newtcolorbox{example}[1]{
    enhanced,
    breakable,
    title={Ejemplo: #1},
    fonttitle=\bfseries,
    colback=green!4,
    colframe=green!40!black,
    boxrule=0.8pt,
    arc=2mm
}

\newtcolorbox{formula}[1]{
    enhanced,
    breakable,
    title={Fórmula / Regla: #1},
    fonttitle=\bfseries,
    colback=purple!4,
    colframe=purple!50!black,
    boxrule=0.8pt,
    arc=2mm
}

\newtcolorbox{exercise}[1]{
    enhanced,
    breakable,
    title={Ejercicio: #1},
    fonttitle=\bfseries,
    colback=gray!5,
    colframe=gray!60!black,
    boxrule=0.8pt,
    arc=2mm
}

\newtcolorbox{note}{
    enhanced,
    breakable,
    title={Nota},
    fonttitle=\bfseries,
    colback=white,
    colframe=gray!50,
    boxrule=0.6pt,
    arc=2mm
}

% ============================================================
% CONFIGURACIÓN FINAL
% ============================================================
\pagestyle{fancy}
\fancyhf{}

\fancyhead[L]{\nouppercase{\leftmark}}
\fancyhead[R]{\materia}

\fancyfoot[C]{\thepage}

\renewcommand{\headrulewidth}{0.4pt}
\renewcommand{\footrulewidth}{0pt}

\fancypagestyle{plain}{
    \fancyhf{}
    \fancyfoot[C]{\thepage}
    \renewcommand{\headrulewidth}{0pt}
}

% ============================================================
% DOCUMENTO
% ============================================================
\begin{document}

% ------------------------------------------------------------
% PORTADA
% ------------------------------------------------------------
\begin{titlepage}

\thispagestyle{empty}

\begin{center}

\vspace*{1cm}

{\large \textsc{\universidad}}

\vspace{0.2cm}

{\large \textsc{\facultad}}

\vspace{0.2cm}

{\large \carrera}

\vspace{3cm}

{\Huge \textbf{\materia}}

\vspace{1cm}

{\LARGE \titulo}

\vspace{0.8cm}

\textit{Comprender es el primer paso para convertir el estudio en conocimiento.}

\vspace{1.2cm}

\rule{0.70\textwidth}{0.5pt}

\vspace{0.5cm}

{\large Apuntes de estudio}

\vspace{0.5cm}

\rule{0.70\textwidth}{0.5pt}

\vfill

{\large \autor}

\vspace{0.3cm}

{\large \anio}

\end{center}

\vfill

\begin{flushleft}
\textit{Este es un modesto aporte a los estudiantes de ``Práctico Profesional'' no pretende sustituir el material oficial de la materia.}
\end{flushleft}

\end{titlepage}

% ------------------------------------------------------------
% ÍNDICE
% ------------------------------------------------------------
\tableofcontents

% ------------------------------------------------------------
% PRESENTACIÓN
% ------------------------------------------------------------
\chapter*{Presentación}
\markboth{Presentación}{Presentación}
\addcontentsline{toc}{chapter}{Presentación}

Estos apuntes reúnen los contenidos de una clase de \textbf{Práctico Profesional} dedicada a mostrar cómo se ejerce la abogacía en la práctica real. El recorrido parte de la vocación de servicio y del cumplimiento de las obligaciones alimentarias; continúa con el peso de lo procesal y, a partir de dos casos prácticos ---un acuse de negligencia en la prueba pericial y un pedido de caducidad de la prueba informativa---, desarrolla la redacción completa de un escrito: título, encabezamiento, objeto, negativa, fundamentos y petitorio.

Además se estudian las costas (arts. 68 y 69 del Código Procesal), la diferencia entre incidente e incidencia, los honorarios y su prescripción, el estilo forense, la relación con el juzgado y las estrategias frente a una providencia: revocatoria, recusación, apelación e instancias supranacionales.

\section*{Utilidad profesional de los contenidos}

Los contenidos de esta clase sirven, en la vida profesional, para:

\begin{itemize}
    \item \textbf{Defender bien a un cliente:} un acuse de negligencia o de caducidad que no se contesta a tiempo puede costar una prueba y, con ella, el juicio.
    \item \textbf{Escribir con claridad y estrategia:} un título que identifica la petición, argumentos que realmente aportan y un petitorio con secuencia lógica.
    \item \textbf{Cuidar el ingreso profesional:} saber quién paga las costas, cuándo se regulan los honorarios y cuándo prescriben.
    \item \textbf{Elegir el camino correcto:} contestar, pedir revocatoria, recusar o apelar, según el caso.
    \item \textbf{Moverse en el juzgado:} conocer al despachante, pedir el alta de patrocinio y asegurar el acceso al expediente.
\end{itemize}

Su utilidad no se limita al litigio. Aprender a argumentar con base normativa y a distinguir un argumento neutro de uno peligroso mejora cualquier negociación o discusión. Gestionar plazos y tiempos es una habilidad útil en cualquier trabajo. La actitud de servicio, en lugar de aprovecharse del conflicto, construye reputación y confianza. Y entender cómo se resuelve un conflicto (el acuerdo primero, la sanción como último recurso) ayuda con los problemas familiares y personales sobre los que tarde o temprano se consulta a quien estudia Derecho: <<¿qué te parece?>>.

\begin{note}
\textbf{Cómo se relacionan los temas: un partido de fútbol.} El apunte propone pensar la clase como un partido.
\begin{itemize}
    \item El \textbf{Bloque 1} es la razón por la que se juega: la vocación de servicio y el <<juego limpio>>.
    \item El \textbf{Bloque 2} es el reglamento y el entrenamiento: sin conocer las reglas (derecho procesal) ni haber practicado (práctica forense) no se puede jugar.
    \item El \textbf{Bloque 3} es armar la jugada de defensa cuando se cobra una falta (el acuse de negligencia).
    \item El \textbf{Bloque 4} es quién paga la multa después de cada falta (costas e incidencias).
    \item El \textbf{Bloque 5} es conocer la cancha y al árbitro (juzgado, despachante, acceso al expediente).
    \item El \textbf{Bloque 6} es la jugada dudosa y el VAR: la caducidad, cómo se comunica la decisión y qué reclamos pueden hacerse (revocatoria, recusación, apelación e instancias superiores).
\end{itemize}
\end{note}

% ============================================================
% CAPÍTULO 1
% ============================================================
\chapter{La profesión como servicio}

\section{Vocación de servicio: no aprovecharse de los conflictos}

El ejercicio de la abogacía consiste en \textbf{brindar un servicio} y una solución, y no en aprovecharse de los conflictos. Lo que prevalece no es una mirada meramente lucrativa: aunque se pretende vivir de brindar ese servicio, este aspecto no debe perderse de vista.

Quienes cursan los últimos años de la carrera suelen recibir consultas de amistades y familiares (<<tengo este problema, ¿qué te parece?>>). Aunque todavía no responden como profesionales, ya tienen un perfil y una trayectoria, y una idea bastante clara de cómo funciona el sistema.

\section{Alimentos: buscar el cumplimiento, no solo el castigo}

\subsection{El cumplimiento como objetivo y el registro de deudores}

En materia de alimentos se buscan mecanismos que promuevan el \textbf{cumplimiento}; el enfoque no es meramente punitivo. La sanción, en muchos casos, busca de algún modo el cumplimiento, y resulta pertinente cuando no queda otra alternativa porque se han agotado otras herramientas. El \textbf{registro de deudores} es un mecanismo adicional que promueve el cumplimiento y lleva al sancionado a cumplir. Con los años, estos mecanismos se van perfeccionando y se buscan nuevas alternativas.

\subsection{La composición y la recalcitrancia}

Aunque se prefiera buscar la \textbf{composición}, existe la \textbf{recalcitrancia}: personas con miradas muy cerradas a las que no les entran las razones. Estas situaciones deben analizarse en el caso concreto.

\subsection{El alimentante que <<paga>> pero no logra un acuerdo}

Es frecuente que la persona a la que se le reclaman alimentos afirme que paga, que trabaja todos los meses, que entrega el dinero y que, sin embargo, se le inicia juicio porque se pretende más de lo que corresponde. Hay situaciones que se vuelven difíciles de manejar: el alimentante dice <<estoy contribuyendo>>, pero no logra construir un acuerdo. Ahí hay que preguntarse por qué no lo logra, porque del otro lado puede existir una insatisfacción legítima o que proviene de otra raíz.

\subsection{La interdisciplina en el práctico}

En el práctico se procura trabajar con \textbf{interdisciplina}, porque existen mecanismos más complejos que no pasan exclusivamente por lo material, que es lo que el abogado aborda de manera más directa. Hay otras motivaciones que requieren profundizar; por eso no conviene quedarse con una mirada limitada.

\begin{important}
El derecho de alimentos busca mecanismos para el cumplimiento y no es meramente punitivo: la sanción (por ejemplo, el registro de deudores) interviene cuando se agotaron otras herramientas. Detrás del conflicto material suele haber motivaciones de otra raíz, por eso se trabaja con interdisciplina.
\end{important}

\section{Compromiso judicial, inmediación y registros}

\subsection{Los juzgados desbordados y la falta de compromiso}

Los juzgados argumentan siempre su desborde. Desde la perspectiva expresada en el apunte, sin embargo, \textbf{falta compromiso}: no puede trabajarse hasta la una y media de la tarde, y quienes ocupan cargos judiciales, remunerados por la sociedad, deben asumir una actitud distinta, incluso en la extensión horaria. No se afirma que no trabajen, pero no es adecuado tratar los asuntos diciendo <<pónganse de acuerdo, porque si no vamos a dictar sentencia>>, lo cual se percibe como una amenaza y deja a las partes sin saber cómo se resolverá.

\subsection{La inmediación en los procesos de familia}

La \textbf{inmediación} se declama mucho, sobre todo en familia. El código contiene normas procesales para los procesos de familia, y esa inmediación debería existir también en la práctica.

\subsection{El registro nacional de deudores y la construcción colectiva}

No se niega la importancia del registro de deudores como mecanismo. Se señala, sin embargo, que un registro provincial o de una sola jurisdicción tiene un límite: el incumplidor se traslada a otra jurisdicción (por ejemplo, para sacar una licencia de conducir). De ahí la referencia a un \textbf{registro nacional}. Se trata de una \textbf{construcción colectiva}: no puede resolverse de manera inmediata, por ejemplo mediante un DNU.

% TODO: contenido incompleto o ambiguo en las notas originales (fragmento inaudible en la grabación sobre el cambio de jurisdicción)

% ============================================================
% CAPÍTULO 2
% ============================================================
\chapter{Lo procesal y la práctica forense}

\section{Un fallo reciente: leer antes de opinar y jurisprudencia}

Ante un fallo de gran repercusión, quienes estudian Derecho reciben preguntas de su entorno. La actitud recomendada es \textbf{leer antes de opinar}: si no se accede de inmediato al fallo, puede leerse primero la noticia periodística y luego el texto, que suele publicarse en el sitio de la Corte.

Suele argumentarse, desde la propia Corte, que sus fallos son decisiones particulares que no pueden extrapolarse a una situación general. Sin embargo, lo que dice la Corte \textbf{marca una línea} e influye en el conjunto; por lo tanto, constituye \textbf{jurisprudencia}, aunque no sea una decisión general.

El fallo comentado muestra además cómo \textbf{lo procesal marca todo}: se planteó una problemática de legitimación por razón colectiva, y las lecturas de lo procesal impactan en lo general. Sin perder de vista el derecho de fondo, los juicios muchas veces se definen por lo procesal.

% TODO: contenido incompleto o ambiguo en las notas originales (no se identifica el fallo comentado)

\section{Derecho procesal y práctica forense: sin ellos no se litiga}

\subsection{Un ejemplo de lo procesal como factor determinante}

\begin{example}{dos procesos y cosa juzgada}
En un juicio en el que intervenía el docente, un colega informó que el resultado había sido adverso para sus clientes. Según se supo, la causa estaba vinculada a una cuestión procesal: había dos procesos y se sostenía que la cosa juzgada de uno influía en el otro. Quedó pendiente analizar si ese planteo era correcto, pero el caso muestra que lo procesal está siempre presente y que por eso es necesario saber derecho procesal.
\end{example}

\subsection{Saber derecho procesal para poder litigar}

Para poder litigar hay que saber \textbf{derecho procesal}. Quien sea un maestro del derecho de fondo en un área, pero no sepa articularlo en el proceso donde esas cuestiones se dirimen, no tendrá buenos resultados.

\subsection{La práctica forense como anexo de lo procesal}

Anexa a lo procesal se encuentra la \textbf{práctica profesional o práctica forense}. Si no se maneja la práctica forense, tampoco se sabe articular el derecho procesal. Es un nivel todavía más concreto y específico.

\begin{important}
Lo procesal es determinante: muchos juicios se definen por una cuestión procesal. El derecho de fondo solo se hace valer si se sabe articular el derecho procesal, y este, a su vez, requiere manejar la práctica forense.
\end{important}

\subsection{Los casos de trabajo y la práctica de redacción}

Se trabajaron dos casos prácticos: uno de \textbf{negligencia} (prueba pericial) y otro de \textbf{caducidad} (prueba informativa). El docente prioriza que los alumnos trabajen siempre un mínimo de redacción, aspecto que suele quedar relegado. Se prevé profundizarlo con el alegato y, en los meses restantes de cursada, con la expresión de agravios.

% ============================================================
% CAPÍTULO 3
% ============================================================
\chapter{Caso 1: contestar un acuse de negligencia}

\section{El caso 1 y su consigna}

\begin{exercise}{Caso 1 --- acuse de negligencia en la prueba pericial}
Usted es abogado de la parte actora en un juicio de daños y perjuicios por daños en un vehículo como consecuencia de una colisión. El día 21 de marzo del 26 se lleva a cabo la audiencia del artículo 360 del Código Procesal y se designa audiencia para que declaren los cinco testigos que su parte propuso, para el 25 de junio, y se sortea perito mecánico. El 25 de junio del 26 se lleva a cabo la audiencia testimonial y declaran tres testigos, porque dos no han podido ser notificados. Usted pide audiencia y se la fijan para el 14 de agosto del 26. El 4 de agosto del 26, la demandada le acusa la negligencia en la producción de la prueba pericial mecánica; argumenta que el perito, debidamente notificado a su domicilio electrónico, no aceptó el cargo y usted no pidió nada al respecto. Además, el plazo de prueba es de 40 días, artículo 367 del Código Procesal, y se encuentra vencido. El juzgado provee, 4 de agosto del 26: <<Del acuse de negligencia, traslado>>. (Fragmento final inaudible.)

\textbf{Consigna:}
\begin{enumerate}
    \item ¿Cómo y cuándo se notifica el despacho que antecede?
    \item ¿Qué puede hacer usted ante ello y, en su caso, qué plazo tiene?
    \item Confeccione una breve pieza procesal al respecto, simulando encabezamiento, argumentación, petitorio, etc.
\end{enumerate}
\end{exercise}

% TODO: contenido incompleto o ambiguo en las notas originales (fragmento final del enunciado inaudible)

\begin{note}
Normas citadas en el enunciado: el \textbf{art. 360} del Código Procesal Civil y Comercial de la Nación (audiencia preliminar, en la que se designa la audiencia de prueba y se sortea perito) y el \textbf{art. 367} (plazo de prueba de 40 días, según el enunciado). El texto de estos artículos no se leyó en clase; conviene consultar el texto oficial vigente.
\end{note}

Frente a un acuse de negligencia hay dos caminos: \textbf{consentir} (no decir nada) o \textbf{contestar y repeler el acuse}. La buena práctica indica repeler el acuse, contestar y confeccionar una breve pieza procesal. La respuesta a la primera pregunta de la consigna (cómo se notifica la providencia) se desarrolla en el capítulo~6.

\section{El título y el encabezamiento del escrito}

\subsection{El título: <<Contesta>> o <<Se rechace>>}

Todo escrito debe llevar un \textbf{título}, porque es fundamental: apunta a que el juzgado identifique la petición y tenga desde el inicio una idea de adónde se dirige el escrito. En este caso, quien responde en principio \textbf{contesta}; \textbf{rechazar} es una petición. Es una cuestión de estilo: puede titularse \textit{<<Contesta acuse de negligencia. Se rechace.>>}.

\subsection{A quién se dirige: la corrección en el trato}

El escrito se dirige al juez o a la jueza (<<señora jueza>>, <<señor juez>>). Importa el tema del género y la corrección en el trato, pero a veces no se sabe quién ocupa el cargo porque el juzgado cambió, o puede haber un subrogante. Para no asumir un género, puede agregarse una barra con la terminación <<a>> (señor/a juez/a).

\subsection{Cómo se presenta quien escribe y en qué persona}

La presentación se hace habitualmente con el nombre, en el carácter de \textbf{apoderado de la parte actora} (o bien la parte se presenta por derecho propio), con el domicilio procesal y la referencia de los autos (<<en autos\ldots>>). Se prefiere la \textbf{primera persona}, que es la forma más natural de comunicación; si no, parece que se hablara de otros. Cuando contestan varias personas, se adapta la redacción. Hay que usar los tiempos verbales correctos, o bien elegir un tiempo verbal y mantenerlo.

\section{Objeto, <<tiempo y forma>> y estilo forense}

\subsection{El objeto}

Muchos escritos incluyen un acápite titulado \textbf{<<Objeto>>}, para que se vea qué es lo que se viene a hacer; en escritos más breves puede omitirse. Se numera con un acápite (por ejemplo, <<1>> o un número romano). La fórmula usual es <<que vengo en legal tiempo y forma\ldots>>, un latiguillo característico.

\subsection{<<En tiempo>>: el tipo de proceso y el plazo}

La expresión <<en tiempo>> supone que se contesta dentro del plazo, por lo que debe identificarse el \textbf{tipo de proceso}. En este caso es un proceso \textbf{ordinario}; en un sumarísimo el plazo sería de tres días, pero en este caso es de cinco días. La fórmula <<en legal tiempo y forma>> resulta armónica.

% TODO: contenido incompleto o ambiguo en las notas originales (el apunte es confuso sobre si el plazo de cinco días corresponde al proceso ordinario o a otro tipo de proceso)

\subsection{El estilo forense y el <<código de rito>>}

Los juzgados están acostumbrados siempre a una misma fórmula. Se habla del <<código de rito>> porque es algo ritual: siempre se mantiene el mismo estilo, el \textbf{estilo forense}. Cada profesión o disciplina tiene su estilo y se lo reproduce. No significa que todo deba ser igual, sino que se busca un lenguaje común, un lenguaje técnico que permita la comunicación; si se escribe de otra manera, no hay comunicación.

\subsection{<<Forma>>: acordadas y reglamentos}

Los escritos tienen una \textbf{forma}. Las \textbf{acordadas} de la Corte son un ejemplo: quien quiere interponer un recurso extraordinario y no cumple con la acordada (por ejemplo, con el número de renglones) puede ver rechazado el recurso de inmediato. También existen \textbf{reglamentos} (por ejemplo, el Reglamento para la Justicia Nacional) que, aunque parecen letra antigua, a veces se invocan para devolver un escrito. Cada fuero tiene su propio reglamento.

\begin{note}
\textbf{Reglamento para la Justicia Nacional.}
\begin{itemize}
    \item \textbf{Art. 51 (honorarios):} el abogado debe clasificar las tareas que realizó, distinguiendo tareas comunes de tareas particulares.
    \item \textbf{Art. 46 (redacción de los escritos):} regula cómo se redactan y presentan los escritos (por ejemplo, el uso de tinta negra en la época del papel).
\end{itemize}
Referencias dadas por el docente, sin lectura del texto.
\end{note}

La clasificación del art. 51 es muy importante en las \textbf{sucesiones}. En un juicio pueden intervenir varios abogados, y las tareas se clasifican así:

\begin{longtable}{@{}L{0.20\textwidth}L{0.36\textwidth}L{0.34\textwidth}@{}}
\toprule
\textbf{Tipo de tarea} & \textbf{Qué comprende} & \textbf{Contra quién se cobra} \\
\midrule
\endfirsthead
\toprule
\textbf{Tipo de tarea} & \textbf{Qué comprende} & \textbf{Contra quién se cobra} \\
\midrule
\endhead
\bottomrule
\endlastfoot
\textbf{Tareas particulares} & Las que el abogado realiza para su propio cliente. & Contra su cliente. \\
\textbf{Tareas comunes} & Las que corresponden a la masa hereditaria, al acervo, y benefician a todos los herederos. & Contra herederos que no fueron sus clientes, que se favorecieron con ese trabajo. \\
\end{longtable}

\begin{example}{tareas comunes en una sucesión}
Si el abogado de un heredero presentó la demanda y luego fueron citados otros herederos, estos se favorecieron con ese trabajo y deben pagarlo aunque no hayan sido sus clientes. Lo mismo ocurre si pidió la declaratoria de herederos. Por eso algunos abogados presentan la demanda aunque los otros herederos aún no aparezcan. Pueden existir sucesiones presentadas por dos abogados que terminan acumulándose en la primera.
\end{example}

\subsection{La firma de los escritos}

Según el reglamento, los escritos debían presentarse con tinta negra, lo que tenía sentido en la época del papel, cuando había que firmar a mano. Hoy la tinta negra es la más fácil de falsificar (mediante escáner, imagen o fotocopia). La firma \textbf{hológrafa} es válida mientras sea efectivamente la firma de la persona, salvo que exista firma electrónica.

Las actas de mediación son hoy digitales, y las partes intervinientes deben estar registradas en el sistema (antes MEPRE; hoy, portal de mediación).

% TODO: contenido incompleto o ambiguo en las notas originales (el nombre del sistema de mediación se menciona de manera poco clara)

\begin{longtable}{@{}L{0.18\textwidth}L{0.34\textwidth}L{0.38\textwidth}@{}}
\toprule
\textbf{Aspecto} & \textbf{Firma electrónica} & \textbf{Firma digital} \\
\midrule
\endfirsthead
\toprule
\textbf{Aspecto} & \textbf{Firma electrónica} & \textbf{Firma digital} \\
\midrule
\endhead
\bottomrule
\endlastfoot
\textbf{Dónde se usa (según el docente)} & En tribunales de la Ciudad de Buenos Aires (Capital). & En la provincia de Buenos Aires. \\
\textbf{Cómo se obtiene} & Mecanismo de firma del sistema de tribunales. & Mediante un token (como un pendrive) con clave. \\
\textbf{Cómo se firma un documento} & --- & Abrir el documento con el programa Adobe, elegir <<firmar con certificado>> e ingresar la contraseña del token. \\
\end{longtable}

Estas herramientas exigen la actitud de aprenderlas y convivir con ellas.

\subsection{Cierre del objeto: <<a contestar el acuse>>}

El objeto se cierra indicando la forma: no se escribe <<acusa de negligencia>>, sino \textit{<<que vengo a contestar el acuse>>}. Pueden usarse fórmulas como <<a tenor de lo siguiente>>, <<a mérito de las razones>> o <<en atención a las consideraciones de hecho y de derecho>>, que suenan más imponentes que <<seguidamente paso a exponer>>. Estas fórmulas asocian con el contexto y hacen ver de qué está hablando el profesional. A continuación se desarrolla el punto en el que <<se define la batalla>>: la negativa y los fundamentos.

\section{La negativa y los fundamentos: análisis crítico de cada argumento}

\subsection{La negativa}

\begin{example}{redacción de la negativa propuesta por un alumno}
\begin{quote}
<<Niego categóricamente que esta parte haya incurrido en inactividad procesal, negligencia o falta de impulso tendiente a la producción de la prueba pericial mecánica ofrecida por esta parte, conforme se demostrará seguidamente>>.
\end{quote}
\end{example}

\subsection{Los fundamentos propuestos}

\begin{example}{fundamentos redactados por un alumno}
\begin{enumerate}
    \item Que la prueba cuya negligencia se acusa fue ofrecida oportunamente en el escrito de demanda, en tiempo y forma.
    \item Contrariamente a lo manifestado de contrario, esta parte realizó todas las diligencias necesarias para lograr la producción de la medida, tal como surge de las actuaciones previas.
    \item La demora invocada es ajena a la conducta procesal de esta parte y responde exclusivamente a la constitución del Tribunal [sic].
    \item La jurisprudencia y doctrina aplicable establecen que el acuse de negligencia tiene carácter restrictivo; no puede prosperar cuando la parte ha demostrado interés en impulsar el procedimiento, evitando así sanciones desproporcionadas que afecten el derecho de defensa en juicio.
    \item Que la ley exige que el juez designe al perito en cuestión, dada la no aceptación del cargo del anterior designado, conforme el artículo 469 del Código Procesal Civil y Comercial de la Nación.
\end{enumerate}
El escrito continúa con el petitorio, como punto cuarto.
\end{example}

\begin{note}
\textbf{Art. 469 del Código Procesal Civil y Comercial de la Nación.} Norma invocada por el alumno: la ley exige que el juez designe al perito ante la no aceptación del cargo del anterior designado. El propio escrito la presenta en el título como <<nombramiento de peritos>>. El docente no leyó el texto del artículo; se consigna la referencia tal como se la invocó en clase.
\end{note}

\subsection{Valoración de cada argumento}

El docente propone ser preciso y valorar cada argumento:

\begin{longtable}{@{}L{0.40\textwidth}L{0.50\textwidth}@{}}
\toprule
\textbf{Argumento del escrito} & \textbf{Valoración} \\
\midrule
\endfirsthead
\toprule
\textbf{Argumento del escrito} & \textbf{Valoración} \\
\midrule
\endhead
\bottomrule
\endlastfoot
\textbf{1. La prueba fue ofrecida oportunamente, en la demanda.} & \textbf{Argumento neutro:} nadie lo cuestiona; <<no aporta>>. \\
\textbf{2. La parte realizó todas las diligencias necesarias.} & \textbf{Argumento peligroso:} en el caso no se ve que la parte haya hecho nada respecto del perito; la parte puede y debe informar al juez y empujar la prueba. \\
\textbf{3. La demora es ajena a la conducta de la parte.} & Se discute junto con el anterior; la jurisprudencia dice que aun así no se exime a la parte cuando el órgano <<se queda dormido>>. \\
\textbf{4. El acuse de negligencia tiene carácter restrictivo.} & \textbf{Argumento fuerte:} la amplitud de la prueba tiene raíz en el derecho de defensa en juicio; la negligencia es su contracara y, como toda pérdida de derechos, se interpreta con criterio restrictivo. \\
\textbf{5. Art. 469: el juez debe designar nuevo perito.} & Argumento con apoyo normativo. \\
\textbf{6. (Agregado en clase) El plazo de prueba no está vencido y hay prueba pendiente.} & Otro argumento interesante, que se suma al escrito. \\
\end{longtable}

El segundo argumento es peligroso porque, aunque sea carga o deber del juez designar al perito, quien es actor está interesado en que la prueba se produzca y puede informar al juez que el perito no aceptó. Un abogado enfocado en la prueba testimonial puede descuidar la pericial, pero las pruebas deben marchar en paralelo.

Respecto del art. 469, existe jurisprudencia según la cual la designación por el juez \textbf{no exime a la parte} cuando <<el órgano se queda dormido>>. El docente aclaró que esta línea argumental no es de su invención: la planteó en dos oportunidades, con resultado favorable en una y desfavorable en otra, y en esta última se le señaló que debió haber impulsado el trámite.

% TODO: contenido incompleto o ambiguo en las notas originales (no queda claro a qué línea argumental se refiere el docente al decir que la ganó y la perdió)

\subsection{Ejemplo: la caducidad de la segunda instancia}

\begin{example}{caducidad de la segunda instancia}
La caducidad de la segunda instancia es difícil de configurar. Opera cuando un recurso debe tratarse y no se trata por una inacción \textbf{de la parte} y no del tribunal. Generalmente se eleva el expediente, la Cámara lo recibe y los autos pasan al acuerdo, pero a veces las partes deben hacer algo: por ejemplo, en el recurso libre, cuando el expediente se pone a los fines de presentar la expresión de agravios. Como rige el principio dispositivo en los procesos civiles y comerciales, a la parte le interesa enviar la cédula. Esto también ocurre actualmente en el fuero laboral, por la ley 27.812.

Ocurre que no se envía la cédula porque el abogado piensa que la mandará el tribunal, y porque la profesión exige hacer muchas cosas a la vez (contestar un telegrama, hacer un escrito, atender a una persona). Por eso hay que saber gestionar el tiempo.
\end{example}

% TODO: contenido incompleto o ambiguo en las notas originales (referencia a la ley 27.812 dicha de manera no verificada)

\subsection{El carácter restrictivo de la negligencia y el proceso que no está a merced de la desidia}

Un alumno señaló que, frente a la jurisprudencia según la cual <<el órgano duerme>>, el carácter restrictivo de la negligencia sería más fuerte. El docente lo valoró: existe un principio de \textbf{amplitud de la prueba}, que tiene su raíz en el \textbf{derecho de defensa en juicio}; la negligencia es su contracara y constituye una restricción. Así como la caducidad de instancia debe verse con efectos restrictivos, la pérdida de derechos se ve siempre con criterio restrictivo. No debe perderse un derecho por un subterfugio procesal.

Pero también ocurre que \textbf{el proceso no está a merced de la desidia}: la negligencia se apoya muchas veces en la desidia y el desinterés. Los juicios tienen sus tiempos y estos no existen para que los juicios se hagan eternos.

\begin{example}{prescripción de la acción penal}
Para ilustrar la idea de que el proceso no está a merced de la desidia, el docente comentó un caso conocido en Tierra del Fuego en el que una persona imputada de homicidio, y que había estado prófuga, pidió la prescripción de la acción penal porque había transcurrido el máximo de la pena. Se trataba de un homicidio agravado por el vínculo; como no existía el femicidio en ese momento, la pena posible no era tan grave. El docente recordó que el homicidio tiene pena de 8 a 25 años, salvo los supuestos de prisión perpetua, y que la prescripción de la acción penal tiene un mínimo y un máximo. Comentó también un argumento del abogado del prófugo: que su cliente <<ya cumplió la pena>> porque tuvo que sustraerse a la justicia durante más de 20 años. El docente señaló que ese argumento no discutía el fondo del asunto.
\end{example}

% TODO: contenido incompleto o ambiguo en las notas originales (el docente duda de los plazos y de las penas mencionadas; el nombre de un caso anterior proviene de la transcripción automática y no es verificable)

\section{El plazo de prueba es ordenatorio y hay prueba pendiente}

\subsection{Hay prueba pendiente: el período de prueba sigue abierto}

Un alumno propuso invocar que existe prueba pendiente de producción: la demandada presentó el acuse el 4 de agosto y la audiencia está fijada para el 14 de agosto, por lo que el período de prueba todavía está abierto. El docente lo consideró otro argumento interesante, y se preguntó si puede analizarse una prueba desconectada de otra, ya que cada una tiene objetos específicos: la prueba pericial responde a los puntos de pericia, y la testimonial interroga a quienes vieron los hechos.

\subsection{El artículo 333: indicar sobre qué hechos declarará cada testigo}

El \textbf{art. 333} del código obliga a la parte a indicar, al proponer a un testigo, respecto de qué hechos va a declarar. Esa norma no prevé una sanción si no se cumple, a diferencia de otras cargas cuyo incumplimiento acarrea la inadmisibilidad de la prueba; aun así, es de buena práctica cumplirla.

\begin{example}{precisión sobre los hechos del testigo}
En una causa por incumplimiento contractual (contrato de obra o de servicio, con una prestación que una parte dice haber realizado y la otra no, y una suerte de excepción de incumplimiento), el juzgado exigió subsanar la indicación genérica de que los testigos declararían <<sobre los hechos de la litis>> o <<sobre que se cumplió el contrato y que hubo incumplimientos de la parte demandada>>. El juzgado quería mayor precisión: qué hechos declararía cada testigo.
\end{example}

\subsection{El plazo de 40 días: perentorio u ordenatorio}

La contraparte alega que el plazo de prueba está vencido. Se trata de un plazo procesal de \textbf{40 días hábiles}. Pero pueden existir pruebas que el propio juzgado fijó, por razones de agenda, más allá de los 40 días. Por eso el plazo es de \textbf{orden}, \textbf{ordenatorio}, y no fatal.

\begin{definition}{plazo ordenatorio}
Plazo cuya finalidad es que los procesos se sustancien en plazos normados. No es un plazo fatal: no clausura automáticamente el período de prueba al vencer, y la prueba que el juzgado fijó más allá de los 40 días se produce igual.
\end{definition}

\begin{longtable}{@{}L{0.17\textwidth}L{0.36\textwidth}L{0.37\textwidth}@{}}
\toprule
\textbf{Criterio} & \textbf{Plazo ordenatorio (de orden)} & \textbf{Plazo fatal} \\
\midrule
\endfirsthead
\toprule
\textbf{Criterio} & \textbf{Plazo ordenatorio (de orden)} & \textbf{Plazo fatal} \\
\midrule
\endhead
\bottomrule
\endlastfoot
\textbf{Finalidad} & Ordenar el proceso para que se sustancie en plazos normados. & Clausurar la posibilidad de actuar al vencer el plazo. \\
\textbf{Cómputo} & Procesal: días hábiles (los 40 días del período de prueba). & --- \\
\textbf{Efecto al vencer} & No clausura automáticamente la prueba: la que el juzgado fijó más allá de los 40 días se produce igual. & Al llegar el día 41 ya no se producirían las medidas y <<lo que está probado está probado>> (el docente descarta que sea el caso). \\
\textbf{Uso en el escrito} & Argumento: <<el plazo de prueba no está vencido>>. & --- \\
\end{longtable}

\begin{important}
El plazo de prueba de 40 días es ordenatorio, no fatal. Si hay prueba pendiente fijada por el propio juzgado, el período probatorio sigue abierto y el acuse de negligencia no puede prosperar por vencimiento del plazo.
\end{important}

\section{Redacción colectiva y petitorio del escrito}

\subsection{Trabajo en equipo y producción colectiva}

El docente recomienda la \textbf{producción colectiva} de escritos. Es desafiante redactar en conjunto, y uno suele poner más la pluma que otro, pero el trabajo en equipo enriquece la idea porque se comparte. Se sugiere compartir los escritos con el grupo, como ocurría en la época en que se subían al campus.

\subsection{Redacción del argumento agregado}

El argumento incorporado en clase puede redactarse así: \textit{<<El plazo no está vencido porque hay prueba pendiente de producción que, según proveído del 25 de junio, nunca fue objetado. La etapa probatoria no está terminada, por lo que el acuse no puede prosperar>>}. Es una versión simplificada de lo propuesto por los alumnos, que suma un argumento más al escrito.

\subsection{El petitorio y su secuencia}

\begin{example}{petitorio redactado por un alumno}
\begin{quote}
<<Cuatro. Petitorio. Por lo expuesto, a vuestra señoría solicito: se tenga por contestado en tiempo y forma el traslado conferido; se rechace el acuse de negligencia articulado por la contraria, manteniendo vigente la prueba ofrecida por esta parte; y se solicite al señor juez la carga legal de designar nuevo perito, conforme lo dispuesto por el artículo 469 del Código Procesal [Civil y Comercial de la Nación]>>.
\end{quote}
\end{example}

El petitorio tiene una \textbf{secuencia}:

\begin{enumerate}
    \item \textbf{Tener por contestado el traslado.} Muchas veces el juzgado dice <<por contestado el traslado>> y se queda ahí; también puede decir <<autos a resolver>>. Primero se tiene por contestado y, una vez que se consiente ese auto, pasan los autos a resolver.
    \item \textbf{Rechazar el acuse.}
    \item \textbf{Mantener vigente la prueba.}
    \item \textbf{Pedir el nuevo sorteo o designación de perito.} Debe hacerse sin <<tirarse contra el juez>> ni recordarle su carga de manera confrontativa, porque no pueden imputarse a la parte omisiones del juzgado. El docente propone agregar: <<solicitamos nuevo sorteo de perito>>. Un alumno señaló que el escrito ya lo incluye en el título (<<Solicita nombramiento de peritos>>), como un agregado.
    \item \textbf{Pedir las costas}, con la fórmula <<las costas a la contraria vencida>>, y un plazo para que el nuevo perito acepte el cargo.
\end{enumerate}

\subsection{Redactar con proporción: incidencia versus demanda}

No se trata de escribir una novela. El escrito analizado es una incidencia, y otros escritos, como la demanda, exigen argumentar con mayor profundidad la pretensión, lo que requiere un esfuerzo intelectual mayor. No es un ensayo: es un escrito concreto por el que se defiende un derecho probatorio que está en crisis.

\begin{important}
El petitorio debe seguir una secuencia: (1) tener por contestado el traslado en tiempo y forma; (2) rechazar el acuse; (3) mantener vigente la prueba; (4) requerir el nuevo sorteo o designación de perito, recordándole al juzgado su carga sin <<tirarse contra el juez>>; y (5) pedir costas a la contraria vencida.
\end{important}

% ============================================================
% CAPÍTULO 4
% ============================================================
\chapter{Costas e incidencias}

\section{Las costas: ¿hay que pedirlas?}

Los alumnos discutieron si corresponde pedir las costas. Una posición sostuvo que no es necesario, porque no se trata de un recurso sino de un acuse de la parte y no habría imposición de costas. Otra sostuvo que sí corresponde pedirlas, porque alguien va a ganar y alguien va a perder. El docente remarcó que se está siempre ligado al éxito y a la derrota, y que conviene conectar el planteo con una \textbf{base normativa}.

\subsection{La base normativa: artículos 68 y 69}

\begin{formula}{art. 68 del Código Procesal Civil y Comercial de la Nación (principio general de costas)}
La parte vencida en el juicio deberá pagar todos los gastos de la contraria, aun cuando ésta no lo hubiese solicitado. Sin embargo, el juez podrá eximir total o parcialmente de esta responsabilidad al litigante vencido, siempre que encontrare mérito para ello, expresándolo en su pronunciamiento, bajo pena de nulidad.
\end{formula}

\begin{formula}{art. 69 del Código Procesal Civil y Comercial de la Nación (costas en los incidentes)}
En los incidentes también regirá lo establecido en el artículo anterior. No se sustanciarán nuevos incidentes promovidos por quien no hubiere sido condenado al pago de las costas en otro anterior, mientras no satisfaga su importe o, en su caso, lo deposite a embargo. Estarán sujetas a este requisito de admisibilidad las incidencias promovidas en el curso de las audiencias, y toda apelación sobre imposición de costas y regulación de honorarios se concederá en efecto diferido, salvo cuando el expediente deba ser remitido a la Cámara como consecuencia del recurso deducido por alguna de las partes contra la resolución que decidió el incidente.
\end{formula}

\begin{note}
Textos leídos en clase por un alumno, con la ortografía corregida; se recomienda cotejarlos con el texto oficial vigente.
\end{note}

Aunque no se soliciten, el juez aplicará el principio igualmente, también en los incidentes, porque el art. 69 remite a lo dispuesto en el artículo anterior.

\begin{important}
Rige el principio objetivo de la derrota (art. 68): la parte vencida paga las costas de la contraria aunque esta no lo haya solicitado, y el juez solo puede eximirla total o parcialmente si encuentra mérito y lo expresa en su pronunciamiento, bajo pena de nulidad. El art. 69 extiende esa regla a los incidentes.
\end{important}

\subsection{El derecho de defensa y la Convención Americana}

Un alumno propuso invocar, en relación con el derecho de defensa, el art. 8 de la Convención Americana sobre Derechos Humanos (garantías judiciales). El docente lo consideró válido, porque se vincula con el derecho probatorio, aunque advirtió que hay que tener cuidado porque ciertas cuestiones de las convenciones se relacionan fundamentalmente con el derecho penal y debe analizarse cómo se vinculan con el derecho en general. El alumno señaló que la Convención va más allá, y el docente coincidió en que se relaciona con el derecho de defensa en sentido amplio.

\begin{note}
\textbf{Convención Americana sobre Derechos Humanos, art. 8.1.} <<Toda persona tiene derecho a ser oída, con las debidas garantías y dentro de un plazo razonable, por un juez o tribunal competente, independiente e imparcial, establecido con anterioridad por la ley, en la sustanciación de cualquier acusación penal formulada contra ella, o para la determinación de sus derechos y obligaciones de orden civil, laboral, fiscal o de cualquier otro carácter>>. El alumno mencionó <<el artículo 8, creo, 8.8>>; el texto transcripto corresponde al art. 8.1, que es el que se vincula con el derecho de defensa. No fue leído en clase.
\end{note}

% TODO: contenido incompleto o ambiguo en las notas originales (el alumno menciona <<8.8>>)

\section{Incidentes e incidencias: un <<microconflicto dentro del conflicto madre>>}

Como el planteo analizado no tramita en un expediente separado, en este caso hay costas en una \textbf{incidencia}. Muchos estudiantes creen que <<incidente>> es cualquier cuestión litigiosa, cuando el término suele reservarse para un conflicto que no es el central y que tramita con un expediente específico.

\begin{definition}{incidente}
Problemática o conflicto que no es el conflicto central del proceso, que tramita en un expediente específico, paralelo al principal, y que existe en función de este. El código contiene un capítulo sobre incidentes (el docente cita el art. 180). Ejemplos: el beneficio de litigar sin gastos (que busca acreditar la imposibilidad de pagar los gastos del proceso principal), la redargución de falsedad, las excepciones previas, los incidentes de nulidad y el incidente de administración en una sucesión.
\end{definition}

\begin{definition}{incidencia}
Contingencia que ocurre dentro del proceso: un <<microconflicto dentro del conflicto madre>> que hay que decidir concretamente para que el proceso pueda avanzar, y que se resuelve en el mismo expediente, sin formar un expediente separado. Ejemplos: un conflicto referido a la prueba (si se produce o no, cómo se termina la etapa) y los modos anormales de terminación del proceso.
\end{definition}

\begin{longtable}{@{}L{0.16\textwidth}L{0.36\textwidth}L{0.38\textwidth}@{}}
\toprule
\textbf{Criterio} & \textbf{Incidente} & \textbf{Incidencia} \\
\midrule
\endfirsthead
\toprule
\textbf{Criterio} & \textbf{Incidente} & \textbf{Incidencia} \\
\midrule
\endhead
\bottomrule
\endlastfoot
\textbf{Trámite} & Se abre un expediente específico, paralelo al principal (hay un capítulo del código sobre incidentes; el docente cita el art. 180). & Se resuelve dentro del mismo expediente; no se forma un expediente separado. \\
\textbf{Qué es} & Problemática o conflicto que no es el conflicto central, pero que existe en función de un expediente principal. & Contingencia o <<microconflicto>> dentro del conflicto madre, que hay que decidir para que el proceso avance. \\
\textbf{Ejemplos citados} & Beneficio de litigar sin gastos; redargución de falsedad; excepciones previas; incidentes de nulidad; incidente de administración en una sucesión. & Un conflicto referido a la prueba (si se produce o no una prueba, cómo se termina la etapa); modos anormales de terminación del proceso. \\
\textbf{Costas} & Rige el art. 69. & El art. 69 las alcanza como <<incidencias>>; puede haber normas más específicas. \\
\end{longtable}

% TODO: contenido incompleto o ambiguo en las notas originales (el docente cita el art. 180 con duda)

Respecto de las excepciones previas y de los modos anormales de terminación del proceso, pueden existir normas más específicas sobre las costas, que deben consultarse.

\subsection{¿Quién decide cómo se tramita? Ejemplos de incidentes}

El abogado interesado debe solicitar la formación de un incidente, aunque el juez también puede decidirlo de oficio. Los incidentes no se forman porque sí, ya que complejizan la actividad, sino porque son necesarios para mantener el \textbf{buen orden del proceso principal}.

\begin{example}{incidentes paralelos al proceso principal}
En una determinación de la capacidad, el objeto es determinar la capacidad de la persona; pero, en paralelo, existe una problemática patrimonial, por lo que se forma un incidente de administración. En una sucesión, el objeto es determinar quiénes son los herederos y qué corresponde a cada uno del acervo del causante; en paralelo se designa un administrador, y todo lo que este hace va a un incidente.
\end{example}

Aplicado al caso, el planteo no es un incidente sino una incidencia: un microconflicto dentro del conflicto madre, que es el derecho que se invoca (en el caso, la indemnización de daños y perjuicios). El conflicto referido a la prueba se resuelve con costas.

\section{Costas, honorarios y su regulación: el <<círculo>> del artículo 69}

La exigencia del art. 69 limita a los litigantes: no pueden promover nuevos incidentes si no pagan las costas de uno anterior, lo que busca \textbf{desalentar la obstrucción} (la <<chicana>>). Pero la regla es <<medio tramposa>>: para pagar las costas deben estar regulados los honorarios, y estos se regulan al final del proceso. Entonces no puede promoverse un incidente porque no se pagaron las costas, pero las costas no están cuantificadas y no pueden pagarse; es una situación que gira sobre sí misma. A veces se fijan las costas y se cuantifica el honorario, pero muchas veces se difiere hasta el final del proceso.

\subsection{Caso real: la sucesión y los honorarios del abogado fallecido}

\begin{example}{sucesión y honorarios de un abogado fallecido}
En una sucesión, el docente y un colega representan a distintos herederos, y habían intervenido antes otros abogados. Al pedir la inscripción de un bien no denunciado originalmente, el juzgado ordenó la inscripción <<previa citación de los profesionales intervinientes>>. Uno de los abogados había fallecido, por lo que hubo que investigar en el Colegio de Abogados y se encontró que tenía herederos y que habían abierto la sucesión. Se les envió un DEO (diligenciamiento electrónico de oficio). Una de las herederas, que trabaja en un juzgado penal, cuestionó por qué no se había planteado la prescripción de los honorarios del abogado fallecido, que inició la sucesión y realizó una \textbf{tarea común}.

El docente explicó que la oportunidad de regulación es al final de la sucesión. El colega indicó que si el abogado cesó en su intervención y pasaron dos años sin que pidiera nada, la inactividad podría dar lugar a la prescripción; planteó que opondría la excepción de prescripción únicamente si los herederos del abogado fallecido piden los honorarios.
\end{example}

% TODO: contenido incompleto o ambiguo en las notas originales (se mencionan dos plazos: cinco años para el honorario regulado y dos años de inactividad)

\begin{longtable}{@{}L{0.38\textwidth}L{0.52\textwidth}@{}}
\toprule
\textbf{Cuestión} & \textbf{Lo que señala el docente} \\
\midrule
\endfirsthead
\toprule
\textbf{Cuestión} & \textbf{Lo que señala el docente} \\
\midrule
\endhead
\bottomrule
\endlastfoot
\textbf{¿Cuándo se regulan los honorarios en una sucesión?} & Al final de la sucesión; no se regulan antes. \\
\textbf{Plazo de prescripción de honorarios regulados} & Cinco años, el máximo del Código Civil y Comercial hoy en día. \\
\textbf{Cese de la intervención del abogado} & Si el abogado cesó y pasaron dos años sin pedir nada, la inactividad podría dar lugar a la prescripción. \\
\textbf{Tarea común} & Quien inicia la sucesión y abre la puerta para que se declaren bienes realiza una tarea común, en beneficio de todos los herederos. \\
\textbf{Estrategia del colega} & Oponer la prescripción recién si los herederos del abogado fallecido piden los honorarios. \\
\end{longtable}

\section{Fórmulas de cierre: <<expresa y ejemplar imposición>> y <<provéase de conformidad>>}

Si no se piden las costas, el escrito suena débil: cabe preguntarse por qué no se pidieron. Un alumno sugirió que quizá se debe a inseguridad sobre cómo se resolverá. Se recomienda pedirlas. Los abogados suelen usar redacciones exageradas: <<con expresa imposición de costas>> o, con otra adjetivación, \textit{<<con expresa y ejemplar imposición de costas>>}, que suena ejemplar, como para advertir que no se hagan más planteos de este tipo.

\subsection{La fórmula <<será justicia>>}

La fórmula \textit{<<provéase de conformidad, será justicia>>} no es obligatoria en ningún código, pero es casi un \textbf{rito}: no se ve un escrito que no la tenga. Un escrito que termine sin ella suena raro, aunque nadie lo objete. Algunos prefieren <<es justo proveer de conformidad>>.

% ============================================================
% CAPÍTULO 5
% ============================================================
\chapter{La práctica en los juzgados}

\section{Conocer el juzgado: estilo, despachante y quién lleva la causa}

Antes había mayor conocimiento de las personas que trabajaban en los juzgados porque existía cierta estratificación; hoy hay mucha más movilidad y a veces depende de quién toque. Es importante saber quién es el \textbf{despachante}, quién es el \textbf{oficial} y quién lleva la causa, y muchas veces conviene hablar con quien la lleva.

Además, cada juzgado tiene su propio estilo. Un alumno señaló que no es lo mismo litigar ante un juzgado más conservador en lo que refiere a las formas que ante otro distinto, por lo que hay que saber dónde se está litigando y, según el caso, ser más laxo al redactar.

\section{Acceso al expediente: el alta del patrocinio y el derecho de defensa}

\begin{example}{una alumna del práctico sin acceso al expediente}
En un juzgado civil, una alumna del práctico no podía ver el expediente de una determinación de capacidad porque no le habían autorizado el alta en el sistema. Se acercó a la mesa de entradas, pidió hablar con quien lleva la causa y un oficial le explicó que debía pedir el \textbf{alta del patrocinio}. La alumna consultó por teléfono con su profesora y resolvió la cuestión.
\end{example}

El abogado a cargo de una comisión del práctico tiene un conjunto de causas del práctico y otro de causas particulares; si no están vinculadas a la comisión, se genera confusión, por lo que es necesario hacer esa vinculación. El práctico tiene mayor presencia en los juzgados de familia que en los patrimoniales.

Se valora la actitud de querer resolver el problema, sin ser <<un pesado>>, sino buscando la respuesta. Es una cuestión formal, pero importante: sin acceso a la causa no hay visibilidad, y el derecho de defensa se ve en riesgo.

\begin{important}
Sin acceso a la causa no hay visibilidad, y sin visibilidad el derecho de defensa corre peligro. Pedir el alta del patrocinio y vincular la causa a la comisión es una cuestión formal, pero decisiva.
\end{important}

% ============================================================
% CAPÍTULO 6
% ============================================================
\chapter{Caso 2: caducidad, notificaciones y estrategia}

\section{El caso 2: caducidad de la prueba informativa}

\begin{exercise}{Caso 2 --- caducidad de la prueba informativa}
Usted es abogado de la parte demandada en un juicio por reivindicación inmobiliaria. El día 21 del 5 del 25 se lleva a cabo la audiencia del artículo 360, se abre la causa a prueba y se ordenan, entre otras medidas, informes al Registro de la Propiedad Inmueble (y a otras dependencias), correspondientes a la prueba de la parte actora. La parte actora libra sus oficios con fecha 25/6/25, y han sido contestados\ldots{} (Fragmento del enunciado inaudible: referido a la prueba informativa.)
\end{exercise}

% TODO: contenido incompleto o ambiguo en las notas originales (el audio del enunciado es fragmentario; se transcribe únicamente lo audible)

En el caso 2 se habla de \textbf{caducidad}, no de negligencia, y el juzgado provee <<traslado>>. Quien pide la caducidad invoca que no se libró el oficio, que no se pidió el reiteratorio o que no se lo pidió dentro de los cinco días.

\section{Negligencia y caducidad: la desidia frente al automatismo}

\begin{definition}{negligencia}
Situación que se produce por la desidia, el desinterés de la parte en la producción de la prueba. Debe ser acusada por la contraria; se corre traslado y la parte puede contestar y repeler el acuse.
\end{definition}

\begin{definition}{caducidad (de la prueba)}
Instituto que opera con cierto \textbf{automatismo}, como consecuencia de una actividad específica que la parte no cumplió. El juzgado tiene que declararla de oficio, por ser un fenómeno procesal de mera comprobación.
\end{definition}

La caducidad tiene raíz normativa. El \textbf{art. 432} se relaciona con la caducidad de la prueba testimonial; los \textbf{arts. 396 y 397} (con duda) se vinculan con la prueba informativa. También se mencionó la caducidad de la prueba confesional. En cuanto a los arts. 34 y 35, se mencionan en relación con las facultades del juez para declararla de oficio.

% TODO: contenido incompleto o ambiguo en las notas originales (numeración de artículos dicha de memoria y con dudas)

Aunque el juzgado deba declararla de oficio, el abogado de la contraria la pide. El juzgado provee <<traslado>> por el \textbf{derecho de defensa}: si se pide la caducidad, hay que escuchar al otro. Ese <<traslado>> se ve en todas las causas. Cabe preguntarse por qué debió contestarse si el juzgado puede declararla de oficio; desde la comprensión de que la caducidad es un fenómeno de mera comprobación, el juzgado tiene que resolver.

\begin{longtable}{@{}L{0.15\textwidth}L{0.36\textwidth}L{0.39\textwidth}@{}}
\toprule
\textbf{Criterio} & \textbf{Negligencia} & \textbf{Caducidad} \\
\midrule
\endfirsthead
\toprule
\textbf{Criterio} & \textbf{Negligencia} & \textbf{Caducidad} \\
\midrule
\endhead
\bottomrule
\endlastfoot
\textbf{Fundamento} & La desidia, el desinterés de la parte en la producción de la prueba. & Un automatismo, producto de una actividad específica que la parte no cumplió en el plazo (por ejemplo, no pedir el reiteratorio dentro de los cinco días). \\
\textbf{Cómo opera} & Debe ser acusada por la contraria; se corre traslado y la parte puede contestar y repeler el acuse. & El juzgado la tiene que declarar de oficio; es un fenómeno procesal de mera comprobación. Aun así, si la pide la contraria, el juzgado corre traslado por el derecho de defensa. \\
\textbf{Interpretación} & Carácter restrictivo: la pérdida de derechos se ve siempre con criterio restrictivo. & Tiene cierto automatismo, pero también se interpreta con efectos restrictivos. \\
\textbf{Ejemplo en clase} & Caso 1: acuse de negligencia por la prueba pericial mecánica. & Caso 2: pedido de caducidad de la prueba informativa (oficios no contestados o reiteratorio fuera de término). \\
\end{longtable}

\section{La providencia <<traslado, notifíquese>>: cómo se notifica y qué se puede hacer}

\subsection{La providencia y la notificación}

La providencia puede formularse así: <<Del pedido de caducidad, traslado. Notifíquese>>. El término \textbf{<<notifíquese>>} significa que hay que hacer una \textbf{cédula}. La pregunta sobre cómo se notificaba en la etapa de prueba remite a los arts. 133 y 135 del código, que conviene internalizar.

\begin{note}
\textbf{Notificaciones.} Art. 133 del Código Procesal Civil y Comercial de la Nación: notificaciones por ministerio de la ley (referencia dada en clase). Art. 135: el docente señala que el <<notifíquese>> siempre remite a este artículo, es decir, a la notificación personal o por cédula. El docente no leyó el texto de los artículos; la numeración fue dicha de memoria.
\end{note}

\subsection{Qué se puede hacer contra la providencia}

Contra una providencia <<traslado, notifíquese>> pueden adoptarse distintas estrategias:

\begin{itemize}
    \item \textbf{Pedir la revocatoria}: <<señor juez, revoque y resuelva directo>>.
    \item \textbf{Mandar la cédula}: es la opción de quien no quiere enfrentarse con el juzgado, pero da al otro la oportunidad de contestar y \textbf{dilata} el trámite.
\end{itemize}

Un alumno agregó un elemento importante: si la providencia está firmada por el \textbf{secretario} y no por el juez, no puede pedirse la revocatoria. Entonces hay que pedirle al juez que revise lo firmado por el secretario; el docente no recordó el nombre del remedio ni los artículos y propuso buscarlos para la próxima clase.

% TODO: contenido incompleto o ambiguo en las notas originales (no se identifica el remedio aplicable a la providencia firmada por el secretario)

\begin{longtable}{@{}L{0.30\textwidth}L{0.29\textwidth}L{0.31\textwidth}@{}}
\toprule
\textbf{Opción frente a <<traslado, notifíquese>>} & \textbf{Ventaja} & \textbf{Riesgo o costo} \\
\midrule
\endfirsthead
\toprule
\textbf{Opción frente a <<traslado, notifíquese>>} & \textbf{Ventaja} & \textbf{Riesgo o costo} \\
\midrule
\endhead
\bottomrule
\endlastfoot
\textbf{Pedir la revocatoria (<<revoque y resuelva directo>>)} & Puede acelerar la resolución sin dar traslado. & Implica <<enfrentarse>> al juzgado; no procede igual si la providencia es del secretario. \\
\textbf{Mandar la cédula} & Evita fricción con el juzgado; es la salida <<más práctica>>. & Le da la oportunidad al otro de contestar: se dilata el trámite. \\
\textbf{Pedir al juez la revisión de lo firmado por el secretario} & Es el remedio que corresponde cuando no firmó el juez (a identificar en el código). & El docente no recordó el artículo: tarea para la próxima clase. \\
\end{longtable}

\begin{note}
\textbf{Tarea pendiente.} Identificar el remedio procesal que corresponde cuando la providencia está firmada por el secretario y no por el juez; el docente propone buscar el artículo del código y trabajarlo en la próxima reunión.
\end{note}

\section{Estrategia procesal: recusación, apelación y la Corte}

\subsection{Enfrentarse al juez cuando no hay alternativa}

Un alumno planteó la conveniencia de no enfrentarse con el juez. El docente respondió que quien quiere ganar el juicio puede tener que enfrentarse al juez en reiteradas ocasiones; plantea la recusación cuando tiene que plantearla, siendo estratégico. No se va al choque porque sí, sino cuando no queda otra alternativa. Se mencionó la \textbf{recusación con causa} y, en el caso de los jueces de Cámara, la posibilidad de plantearla incluso sin causa. La recusación no es aplicable al caso de la providencia analizada, y se la mencionó como ejemplo de actitud estratégica.

El código debe estudiarse para encontrar todas las respuestas de las que uno puede valerse; quien lo maneja bien <<tiene un arma>>.

\subsection{Apelar y las instancias superiores}

Muchas veces corresponde \textbf{apelar}: ante una decisión desfavorable en primera instancia, la primera reacción es apelar y acudir a la Cámara para que escuche, aunque a veces no escuche.

Un alumno preguntó si, ante el fallo comentado, podría acudirse a la Corte. El docente respondió que es difícil. Señaló que la Argentina acepta la \textbf{jurisdicción internacional}: algunos fallos de la Corte han motivado recomendaciones de la Comisión Interamericana, y la Corte Interamericana no cambia el fallo, pero el Estado, para no tener responsabilidad, suele actuar en consecuencia. No es una puerta de revisión. El sistema supranacional es importante para que la Corte nacional no crea que sus decisiones no tendrán consecuencias porque nadie puede revisarlas.

% TODO: contenido incompleto o ambiguo en las notas originales (la descripción del funcionamiento del sistema interamericano es parcial)

\begin{important}
El abogado debe ser estratégico: no se va al choque con el juez porque sí, sino cuando no queda otra alternativa (recusación, apelación). El código hay que estudiarlo para encontrar todos los remedios disponibles, y el sistema supranacional existe para que ningún tribunal, ni siquiera el más alto, crea que sus decisiones no pueden tener consecuencias.
\end{important}

% ============================================================
% CORNELL
% ============================================================
\chapter{Cuadros Cornell de repaso}

El método Cornell organiza el estudio en una columna de preguntas o conceptos clave, una columna de notas o respuestas y una síntesis final. Se presenta un cuadro por cada bloque temático.

\section{Bloque 1: La profesión como servicio}

\begin{longtable}{@{}L{0.26\textwidth}L{0.68\textwidth}@{}}
\toprule
\textbf{Preguntas / palabras clave} & \textbf{Notas / respuestas} \\
\midrule
\endfirsthead
\toprule
\textbf{Preguntas / palabras clave} & \textbf{Notas / respuestas} \\
\midrule
\endhead
\bottomrule
\endlastfoot

\textbf{¿Para qué ejerce el abogado?} &
Para brindar un servicio y una solución, y vivir de ello, sin una mirada puramente lucrativa y sin aprovecharse de los conflictos. \\

\textbf{¿Qué lugar ocupa la sanción en materia de alimentos?} &
Se buscan mecanismos de cumplimiento, no solo punitivos. La sanción (por ejemplo, el registro de deudores) interviene cuando se agotaron otras herramientas y busca llevar al sancionado a cumplir. \\

\textbf{Composición, recalcitrancia e interdisciplina} &
Se prefiere la composición, pero hay personas con miradas cerradas (recalcitrancia). Cuando el alimentante <<paga>> y no logra un acuerdo, suele haber una insatisfacción de otra raíz, por lo que se trabaja con interdisciplina y se analiza el caso concreto. \\

\textbf{¿Qué se señala sobre el compromiso judicial y la inmediación?} &
Los juzgados están desbordados, pero falta compromiso. La inmediación, tan declamada en familia, debería existir en la práctica. \\

\textbf{¿Qué límite tiene un registro de deudores provincial?} &
El incumplidor puede trasladarse a otra jurisdicción. De ahí la idea de un registro nacional, que depende de una construcción colectiva. \\

\midrule
\multicolumn{2}{L{\dimexpr0.94\textwidth+2\tabcolsep\relax}}{%
\textbf{Resumen:} El abogado ejerce para brindar un servicio. En alimentos se busca primero el cumplimiento y el acuerdo, con la sanción como recurso cuando se agotaron otras herramientas; detrás del conflicto material suele haber motivaciones de otra raíz, por lo que se trabaja con interdisciplina.} \\
\bottomrule
\end{longtable}

\section{Bloque 2: Lo procesal y la práctica forense}

\begin{longtable}{@{}L{0.26\textwidth}L{0.68\textwidth}@{}}
\toprule
\textbf{Preguntas / palabras clave} & \textbf{Notas / respuestas} \\
\midrule
\endfirsthead
\toprule
\textbf{Preguntas / palabras clave} & \textbf{Notas / respuestas} \\
\midrule
\endhead
\bottomrule
\endlastfoot

\textbf{¿Por qué un fallo <<particular>> es jurisprudencia?} &
Porque lo que dice la Corte marca una línea e influye en el conjunto, aunque no sea una decisión general. Antes de opinar sobre un fallo hay que leerlo. \\

\textbf{¿Por qué lo procesal es determinante?} &
Porque muchos juicios se definen por lo procesal. El derecho de fondo solo se hace valer si se sabe articularlo en el proceso. \\

\textbf{¿Qué se necesita para poder litigar?} &
Saber derecho procesal y, además, manejar la práctica forense, que es el nivel más concreto y específico de la práctica profesional. \\

\midrule
\multicolumn{2}{L{\dimexpr0.94\textwidth+2\tabcolsep\relax}}{%
\textbf{Resumen:} Lo procesal es determinante. Para litigar hay que dominar el derecho procesal y la práctica forense, además del derecho de fondo; y los fallos de la Corte, aunque sean particulares, marcan una línea y constituyen jurisprudencia.} \\
\bottomrule
\end{longtable}

\section{Bloque 3: Caso 1, contestar un acuse de negligencia}

\begin{longtable}{@{}L{0.26\textwidth}L{0.68\textwidth}@{}}
\toprule
\textbf{Preguntas / palabras clave} & \textbf{Notas / respuestas} \\
\midrule
\endfirsthead
\toprule
\textbf{Preguntas / palabras clave} & \textbf{Notas / respuestas} \\
\midrule
\endhead
\bottomrule
\endlastfoot

\textbf{¿Qué se puede hacer ante un acuse de negligencia?} &
Consentir (no decir nada) o contestar y repeler el acuse. La buena práctica indica contestar con una breve pieza procesal: título, encabezamiento, objeto, negativa, fundamentos y petitorio. \\

\textbf{¿Cómo se titula y encabeza el escrito?} &
El título debe permitir que el juzgado identifique la petición (<<Contesta acuse de negligencia. Se rechace>>). Se cuida el trato (señor juez / señora jueza), la presentación (apoderado o por derecho propio), el domicilio procesal y el uso de la primera persona. \\

\textbf{¿Qué incluye el objeto?} &
La indicación de que se viene <<en legal tiempo y forma>> a contestar el acuse, atendiendo al tipo de proceso y al plazo. \\

\textbf{¿Qué importancia tienen la forma y el estilo forense?} &
Los escritos tienen una forma (acordadas, reglamentos como los arts. 51 y 46 del Reglamento para la Justicia Nacional). El estilo forense busca un lenguaje común. Hay firma electrónica en tribunales de la Ciudad y firma digital con token en la provincia de Buenos Aires. \\

\textbf{¿Qué argumentos sirven para repeler el acuse?} &
Es neutro que la prueba fue ofrecida a tiempo; es peligroso afirmar que se hicieron todas las diligencias si la parte no hizo nada respecto del perito. Son fuertes el carácter restrictivo de la negligencia (la amplitud de la prueba tiene raíz en el derecho de defensa), el art. 469 y que el plazo de prueba no está vencido porque hay prueba pendiente. \\

\textbf{¿El plazo de prueba de 40 días es fatal?} &
No: es un plazo procesal de días hábiles y ordenatorio. Si el juzgado fijó prueba más allá de los 40 días, el período probatorio sigue abierto. El art. 333 exige indicar sobre qué hechos declarará cada testigo (de buena práctica, sin sanción expresa). \\

\textbf{¿Cómo se ordena un petitorio?} &
(1) Tener por contestado el traslado en tiempo y forma; (2) rechazar el acuse; (3) mantener vigente la prueba; (4) pedir nuevo sorteo o designación de perito sin <<tirarse contra el juez>>; (5) pedir costas a la contraria vencida. \\

\midrule
\multicolumn{2}{L{\dimexpr0.94\textwidth+2\tabcolsep\relax}}{%
\textbf{Resumen:} Ante un acuse de negligencia corresponde contestar con un escrito ordenado, con título que identifique la petición, argumentos que realmente aporten (distinguiendo los neutros, los peligrosos y los sólidos) y un petitorio con secuencia. El plazo de prueba es ordenatorio y la pérdida de derechos se interpreta con criterio restrictivo.} \\
\bottomrule
\end{longtable}

\section{Bloque 4: Costas e incidencias}

\begin{longtable}{@{}L{0.26\textwidth}L{0.68\textwidth}@{}}
\toprule
\textbf{Preguntas / palabras clave} & \textbf{Notas / respuestas} \\
\midrule
\endfirsthead
\toprule
\textbf{Preguntas / palabras clave} & \textbf{Notas / respuestas} \\
\midrule
\endhead
\bottomrule
\endlastfoot

\textbf{¿Quién paga las costas? (arts. 68 y 69)} &
La parte vencida paga los gastos de la contraria aunque no se los pida; el juez solo puede eximirla total o parcialmente si encuentra mérito y lo expresa, bajo pena de nulidad. El art. 69 lo extiende a los incidentes. \\

\textbf{¿Qué limitación impone el art. 69?} &
No pueden promoverse nuevos incidentes mientras no se paguen las costas de uno anterior. Es un límite que busca desalentar la chicana, pero resulta <<tramposo>> porque los honorarios suelen regularse recién al final. \\

\textbf{Incidente versus incidencia} &
El incidente tramita en expediente separado y existe en función del principal (beneficio de litigar sin gastos, redargución de falsedad, administración en una sucesión). La incidencia es un <<microconflicto dentro del conflicto madre>> que se resuelve en el mismo expediente para que el proceso avance. \\

\textbf{Honorarios y prescripción} &
Los honorarios se regulan al final del proceso. Las tareas comunes en una sucesión benefician a todos los herederos. El honorario regulado prescribe a los cinco años; la inactividad del abogado puede dar lugar a la prescripción. \\

\textbf{¿Son obligatorias las fórmulas de cierre?} &
No lo son, pero son parte del rito: <<con expresa y ejemplar imposición de costas>> y <<provéase de conformidad, será justicia>>. \\

\midrule
\multicolumn{2}{L{\dimexpr0.94\textwidth+2\tabcolsep\relax}}{%
\textbf{Resumen:} Rige el principio de la derrota (art. 68), extendido a los incidentes (art. 69). Hay que distinguir incidente (expediente separado) de incidencia (microconflicto dentro del mismo expediente) y tener presente que los honorarios se regulan al final y prescriben a los cinco años.} \\
\bottomrule
\end{longtable}

\section{Bloque 5: La práctica en los juzgados}

\begin{longtable}{@{}L{0.26\textwidth}L{0.68\textwidth}@{}}
\toprule
\textbf{Preguntas / palabras clave} & \textbf{Notas / respuestas} \\
\midrule
\endfirsthead
\toprule
\textbf{Preguntas / palabras clave} & \textbf{Notas / respuestas} \\
\midrule
\endhead
\bottomrule
\endlastfoot

\textbf{¿Por qué hay que conocer el juzgado?} &
Porque cada juzgado tiene su estilo y porque es clave saber quién es el despachante, quién es el oficial y quién lleva la causa; hoy hay mucha más movilidad y depende de quién toque. \\

\textbf{¿Qué pasa si no se tiene acceso al expediente?} &
Sin acceso a la causa no hay visibilidad y el derecho de defensa corre peligro. Hay que pedir el alta del patrocinio y vincular la causa a la comisión. \\

\midrule
\multicolumn{2}{L{\dimexpr0.94\textwidth+2\tabcolsep\relax}}{%
\textbf{Resumen:} Conocer al juzgado y asegurar el acceso al expediente son cuestiones formales pero decisivas: de ellas depende la visibilidad de la causa y, con ella, el derecho de defensa.} \\
\bottomrule
\end{longtable}

\section{Bloque 6: Caso 2, caducidad, notificaciones y estrategia}

\begin{longtable}{@{}L{0.26\textwidth}L{0.68\textwidth}@{}}
\toprule
\textbf{Preguntas / palabras clave} & \textbf{Notas / respuestas} \\
\midrule
\endfirsthead
\toprule
\textbf{Preguntas / palabras clave} & \textbf{Notas / respuestas} \\
\midrule
\endhead
\bottomrule
\endlastfoot

\textbf{Negligencia versus caducidad} &
La negligencia responde a la desidia en producir la prueba y debe acusarse. La caducidad opera con cierto automatismo, por una actividad específica incumplida (por ejemplo, no pedir el reiteratorio en cinco días), y el juzgado la declara de oficio, aunque corre traslado por el derecho de defensa. \\

\textbf{<<Traslado, notifíquese>>: ¿qué se hace?} &
<<Notifíquese>> implica cédula (arts. 133 y 135). Se puede pedir revocatoria (<<revoque y resuelva directo>>) o mandar la cédula, que dilata. Si la providencia la firma el secretario, hay que identificar el remedio correspondiente: tarea para la próxima clase. \\

\textbf{¿Cuándo conviene enfrentarse con el juez?} &
No se va al choque porque sí, sino cuando no queda otra alternativa: recusación (con causa; en Cámara también sin causa) o apelación. Hay que estudiar el código para tener todas las respuestas. \\

\textbf{¿Qué función cumple el sistema supranacional?} &
Evita que se crea que las decisiones de la Corte nacional no tienen consecuencias porque nadie puede revisarlas. No es una puerta de revisión del fallo, pero el Estado suele actuar en consecuencia. \\

\midrule
\multicolumn{2}{L{\dimexpr0.94\textwidth+2\tabcolsep\relax}}{%
\textbf{Resumen:} La negligencia se acusa y se contesta; la caducidad opera con automatismo, aunque se corre traslado. Frente a una providencia hay que evaluar revocatoria o cédula, y el abogado debe ser estratégico: conocer el código, recusar o apelar cuando no queda alternativa y recordar la existencia del sistema supranacional.} \\
\bottomrule
\end{longtable}

\section{Síntesis final de la clase}

La clase integró la ética de la profesión con la técnica procesal. El abogado ejerce para brindar un servicio y, ante un conflicto como el de alimentos, busca primero el cumplimiento y el acuerdo, con la sanción como último recurso. Para que ese servicio sea eficaz hay que dominar lo procesal y la práctica forense, porque los juicios muchas veces se definen allí.

A partir de dos casos, se aprendió a contestar un acuse de negligencia con un escrito ordenado (título que identifica la petición, encabezamiento correcto, objeto, negativa, fundamentos que realmente aporten y petitorio con secuencia), distinguiendo argumentos neutros, peligrosos y sólidos, y recordando que el plazo de prueba es ordenatorio y que la pérdida de derechos se interpreta con criterio restrictivo. Se estudiaron las costas (arts. 68 y 69), la diferencia entre incidente e incidencia, los honorarios y su prescripción, así como la diferencia entre negligencia y caducidad. Finalmente, se vio cómo reaccionar ante una providencia (<<traslado, notifíquese>>), cuándo recusar o apelar y por qué importa el sistema supranacional.

El hilo conductor es uno solo: conocer las reglas, escribir con precisión, actuar con estrategia y mantener siempre el compromiso con el problema del cliente y con su derecho de defensa.

% ============================================================
% APÉNDICES
% ============================================================
\appendix

\chapter{Referencias normativas mencionadas en la clase}

Todas las referencias fueron dadas por el docente o por los alumnos durante la clase, en su mayoría de memoria y sin lectura del texto, salvo los arts. 68, 69 y el art. 8.1 de la Convención Americana, cuyos textos se transcribieron en el capítulo~4. Se recomienda confirmar la numeración y el texto en la normativa vigente.

\begin{longtable}{@{}L{0.30\textwidth}L{0.60\textwidth}@{}}
\toprule
\textbf{Norma} & \textbf{Referencia en la clase} \\
\midrule
\endfirsthead
\toprule
\textbf{Norma} & \textbf{Referencia en la clase} \\
\midrule
\endhead
\bottomrule
\endlastfoot
CPCCN, art. 360 & Audiencia preliminar, en la que se designa la audiencia de prueba y se sortea perito (casos 1 y 2). \\
CPCCN, art. 367 & Plazo de prueba de 40 días, según el enunciado del caso 1. \\
CPCCN, art. 333 & Obliga a indicar sobre qué hechos declarará cada testigo; no prevé sanción por incumplimiento. \\
CPCCN, art. 469 & Invocado por un alumno: el juez debe designar perito ante la no aceptación del cargo del anterior designado. \\
CPCCN, arts. 68 y 69 & Principio general de costas y costas en los incidentes. \\
CPCCN, art. 180 & Capítulo sobre incidentes (cita del docente, con duda). \\
CPCCN, arts. 396 y 397 & Relacionados con la caducidad de la prueba informativa (mención poco clara). \\
CPCCN, art. 432 & Caducidad de la prueba testimonial (según el docente). \\
CPCCN, arts. 34 y 35 & Mención parcial, en relación con la declaración de la caducidad de oficio. \\
CPCCN, arts. 133 y 135 & Notificaciones; el <<notifíquese>> remite al art. 135 (notificación personal o por cédula). \\
Reglamento para la Justicia Nacional, art. 51 & Honorarios: clasificar tareas comunes y particulares. \\
Reglamento para la Justicia Nacional, art. 46 & Redacción de los escritos. \\
Convención Americana sobre Derechos Humanos, art. 8 & Garantías judiciales, vinculadas al derecho de defensa. \\
Ley 27.812 & Mencionada en relación con el fuero laboral y el envío de cédula para la expresión de agravios. \\
Código Civil y Comercial & Plazo de prescripción de cinco años para los honorarios regulados (<<el máximo>>, según el docente). \\
\end{longtable}

\end{document}
